%%%                     Súbor ZfSMEscenar.tex                             %%%

% Obsahuje vlastné makrá pre texty divadelných scenárov v knihách zo stránok
% http://zlatyfond.sme.sk, ktoré boli  odladené  pre  hru Chudobná rodina od
% Boženy Slančíkovej-Timravy.  Pre spracovanie iných hier bude asi  potrebné
% doplniť niektoré makrá. Pre vlastné makrá je okrem samotného systému plain
% pdfTeX potrebný aj CSTeX a súbor makier OPmac a makrá ZfSME.tex.



%%% Vlastné makrá (začínajú sa obyčajne veľkým písmenom) %%%

% V prípade že sa nepoužíva TeXonWeb -- https://tex.mendelu.cz pomocou
%\input TeXonWeb.tex % je potrebné definovať makro \ToWp :
\ifx\ToWp\undefined \def\ToWp{} \fi

% Pre staršie verzie CSTeX nemusí byť \elq definované
\ifx\elq\undefined \chardef \elq 96 \fi

% Prednastavené hodnoty, pokiaľ neboli definované v konfigurácii:
\ifx\ZFroot\undefined \def\ZFroot{} \fi % ak nie je iná cesta k ZfSME.tex
\ifx\Vers\undefined \input \ZFroot ZfSME \fi % ak nie je načítaný ZfSME.tex

% Predefinovaie makra \TypoInfo pre použitie v českej aj slovenskej verzii
% textu (pomocou makier \sdef a \mtext z balíka OPmac):
\sdef{mt:typoinfo:sk}%
{Sadzba programom \ulink[https://ctan.org]{\TeX} písmom \FontType.}
\sdef{mt:typoinfo:cs}%
{Sazba programem \ulink[https://ctan.org]{\TeX} písmem \FontType.}
\def\TypoInfo{% Vytlačí info o TeX a použitom fonte
\ifdim\vsize > 110mm \vskip 48pt\noindent \else \medskip\noindent \fi
\mtext{typoinfo}}


% Použitie kapitálok (s tvarom veľkých písmen pre základný rez Roman)
\def\CR{\caps\rm}
\def\CB{\caps\bf}
% Kapitálky nefungujú v písme CM/CS, preto:
\ifx\FontType\CMCS 
\let\CR=\empty 
\let\CB=\bf 
% Nefungujú tu ani ,jednoduché` úvodzovky (viď TPP s. 16 a s. 134):
\def\clq{,}
\def\crq{`}
\fi
% Kapitálky nefungujú ani v písme Charter, preto:
\def\tmp{Charter}
\ifx\FontType\tmp
\let\CR=\empty 
\let\CB=\bf 
\fi

\def\ScenaPopis#1\par{% Popis scény na začiaku dejstva (1 odstavec)
\noindent{\it #1\par}\medskip\noindent}

\def\SP#1){% popis scény v rámci repliky - medzi "(" a ")"
{\it #1)\/}}

\def\SPP#1){\noindent % popis scény medzi replikami - medzi "(" a ")"
{\it #1)\/}\smallskip}


\def\Obsadenie#1\par{% Zoznam účinkujúcich vo výstupe (1 odstavec)
\nobreak\noindent{\it #1}\medskip
}

\def\Rep#1:#2\par{\def\tmp{#1} % Replika v tvare osoba:text repliky
\ifx\tmp\Osoba 
\noindent {\CB #1}{\hangindent=.5\parindent #2\smallskip}%
\else\ifx\tmp\OsobaB
\noindent {\CB #1}{\hangindent=.5\parindent #2\smallskip}%
\else\ifx\tmp\OsobaC
\noindent {\CB #1}{\hangindent=.5\parindent #2\smallskip}%
\else
\noindent {\CR #1}{\hangindent=.5\parindent #2\smallskip}%
\fi\fi\fi}

\def\RepC#1\par{% pokračovanie repliky (napríklad prerušenej veršom)
{\hangindent=.5\parindent \parindent=\hangindent #1\smallskip}
}

\def\Opona#1\par{%
    \removelastskip\nobreak\MedSkip\centerline{\it(Opona #1\kern-.25em.)}
}

% Slovo Konec a jeho formát pre celý riadok na konci dejstva (CZ)
\def\Konec#1\par{\nobreak\MedSkip\noindent{\fC\it Konec\ #1\par}}


%% https://groups.google.com/g/comp.text.tex/c/SbzyM6ZdVwo/m/HuDeL4vKPgsJ
%% (viď Anil Trivedi, 5. 10. 1993, 13:20:49)
\long\def\Ignore#1\endIgnore{\relax}
\def\endIgnore{\relax} % možné použiť priamo \Ignore ... \endIgnore
%% \Ignore je možné použiť aj vo vnútri ďaľšieho makra,  takže funguje aj
\def\ifIgnoreVystup#1{% \ifIgnoreVystup{osoba1,osoba2,...} ... \endIgnore
\ifx\OsobyFilter\undefined \def\OsobyFilter{0}\fi
\ifx\Osoba\undefined \def\Osoba{0}\fi
\ifx\OsobaB\undefined \def\OsobaB{0}\fi
\ifx\OsobaC\undefined \def\OsobaC{0}\fi
\if\Osoba0 \else \ifx^#1^ \else % Ak Osoba=0, alebo je zoznam osôb prázdny
   \ifnum\OsobyFilter=1         % táto vetva sa nevykonáva
          \expandafter\jevzozname \Osoba\in{#1}\ifincluded \else 
          \expandafter\jevzozname \OsobaB\in{#1}\ifincluded \else 
          \expandafter\jevzozname \OsobaC\in{#1}\ifincluded \else 
          \Ignore    % výstup sa netlačí len pri kombinácii ak je nastavený
          \fi\fi\fi  % filter a osoba nie je v zozname pre výstup
   \fi\fi\fi
}
%% TPP, str. 71--72
\newif\ifincluded
\def\jevzozname#1\in#2{% #1 = slovo, #2 = zoznam
\def\tmp##1,#1,##2\end{\ifx^##2^\includedfalse\else\includedtrue\fi}%
\tmp,#2,#1,\end
}

\endinput

%%% Náhrady vo vim (aplikovať až po logickom vyznačení ostatných častí hry
%   (makrá \sec, \secc, seccc, \Perex, \ScenaPopis, \Obsadenie, \Opona, ...
%   okrem makra pre vyznačenie repliky \Rep)):

- pred doplnením vyznačenia replík je teraz vhodné aplikovať makro vim
  pre preformátovanie "rozbitých" odstavcov:
:'<,'>call Justify("tw",3)

- odstavec sa začína dvomi veľkými písmenami, preto by to mala byť replika
(makro \Rep):
:%s/\n\(\n^[A-Z][A-Z]\)/\r\r\\Rep \1/gc

- pre hru Páva obsahovali prvé dva znaky diakritiku (ĎURKO, IĽKA, ŠTEFAN),
  preto aj:
:%s/\n\(\n^[Ď][U]\)/\r\r\\Rep \1/gc
:%s/\n\(\n^[I][Ľ]\)/\r\r\\Rep \1/gc
:%s/\n\(\n^[Š][T]\)/\r\r\\Rep \1/gc
